\section{Experimental Evaluation}
\label{sec:exp}

\stitle{Hardware.}
All algorithms are implemented in C++ and compiled with \texttt{-O3}\footnote{Code and data at \url{https://github.com/anonymous/nsi}.}.
%
Experiments were conducted on a server with $96$ cores and $503$GB of memory, running Linux; all runs are single-threaded and executed one at a time to avoid contention.

\stitle{Algorithms.}
We compare four methods.
%
\specnd{} is our one-build sweep (Algorithm~\ref{alg:specnd}).
%
\specndbase{} is an ablation that shares the build but peels every cell in full, disabling the chain certificate; it isolates the contribution of certification from the shared build.
%
\cnd~\cite{LocalNuclear} is the state-of-the-art per-cell algorithm, run once per cell of the range.
%
\percell{} is our own engine run cell by cell, which serves as a same-engine baseline for the certificate's effect.

\stitle{Datasets.}
We use twelve real graphs from four domains, listed in Table~\ref{tab:datasets}.
%
Web graphs (\daWIT, \daWUK) and structural and fluid-dynamics matrices (\daPKa, \daPWT, \daLDR, \daRAE, \daPKb, \daNAS) have large, tightly overlapping cliques; collaboration graphs (\daDBL, \daAST, \daHEP) are clique-rich but sparser; social graphs (\daEPN, \daYTB) have many small, loosely overlapping cliques.
%
The largest, \daLDR, has over $20$ million edges.

\stitle{Metrics.}
We report end-to-end time and peak memory for computing a whole spectrum, a range of four cells for order $r$ (three for the social graphs).
%
We set a two-hour budget per cell and a $300$GB memory budget; a method that exceeds either on a cell cannot complete that cell, and we mark such cells as out of budget.
%
For per-cell baselines, a spectrum is out of budget when any of its cells is.
%
Correctness is established by the theorems of Sections~\ref{sec:theory}--\ref{sec:algo}; on every cell where a per-cell run completes, the output of \specnd{} is bit-identical to it.
\todo{E2: repeat 3x, report medians; currently single-trial}

% -------- Datasets table --------
\begin{table}[t]
\centering
\caption{Datasets.}
\label{tab:datasets}
\small
\begin{tabular}{llrr}
\toprule
Domain & Graph & $n$ & $m$ \\
\midrule
\multirow{2}{*}{Web}
 & \daWIT{} (web-it-2004)   & 509{,}338   & 7{,}178{,}413 \\
 & \daWUK{} (web-uk-2005)   & 129{,}632   & 11{,}744{,}049 \\
\midrule
\multirow{6}{*}{Matrix}
 & \daRAE{} (raefsky3)      & 21{,}200    & 733{,}783 \\
 & \daPKb{} (pkustk11)      & 87{,}804    & 2{,}565{,}054 \\
 & \daPKa{} (pkustk13)      & 94{,}893    & 3{,}260{,}967 \\
 & \daNAS{} (nasasrb)       & 54{,}870    & 1{,}311{,}227 \\
 & \daPWT{} (pwtk)          & 217{,}891   & 5{,}653{,}221 \\
 & \daLDR{} (ldoor)         & 909{,}537   & 20{,}770{,}807 \\
\midrule
\multirow{3}{*}{Collab.}
 & \daDBL{} (com-dblp)      & 317{,}080   & 1{,}049{,}866 \\
 & \daAST{} (ca-AstroPh)    & 18{,}772    & 198{,}050 \\
 & \daHEP{} (ca-HepPh)      & 12{,}008    & 118{,}489 \\
\midrule
\multirow{2}{*}{Social}
 & \daEPN{} (soc-Epinions)  & 75{,}879    & 405{,}740 \\
 & \daYTB{} (com-youtube)   & 1{,}134{,}890 & 2{,}987{,}624 \\
\bottomrule
\end{tabular}
\end{table}

\expsection{Spectrum Time}
Table~\ref{tab:main} reports the end-to-end time and memory to compute a whole spectrum with \specnd{} and \cnd.
%
\specnd{} outperforms \cnd{} by $20$ times on average and up to $511$ times on \daRAE, where \cnd{} takes $307$ seconds and \specnd{} takes $0.6$ seconds.
%
On five graphs (\daWIT, \daWUK, \daPKa, \daPWT, \daLDR), \cnd{} cannot complete any cell within the $300$GB budget, while \specnd{} finishes the whole spectrum in $7.5$ to $81$ seconds.
%
On \daDBL{} at $r{=}5$, \cnd{} takes $4{,}261$ seconds and \specnd{} takes $28$ seconds, a $151\times$ speedup.
%
On the collaboration graphs, \cnd{}'s per-cell counting is nearly flat in $s$, so the spectrum comparison is close: \specnd{} is $1.7\times$ faster on \daAST{} and $0.8\times$ on \daHEP{}, since there almost every $r$-clique must still be peeled at the boundary cell and the counting cost per cell is small.
%
This experiment demonstrates that \specnd{} wins by orders of magnitude wherever the per-cell state is large, and stays competitive where it is not.

% -------- Main table --------
\begin{table}[t]
\centering
\caption{Whole-spectrum time and memory.}
\label{tab:main}
\small
\begin{tabular}{llrrrr}
\toprule
 & & \multicolumn{2}{c}{\specnd} & \multicolumn{2}{c}{\cnd} \\
\cmidrule(lr){3-4}\cmidrule(lr){5-6}
Graph & $(r,s)$ & time & mem & time & mem \\
\midrule
\daWIT  & $(3,4$--$7)$ & 8.1\,s  & 0.5\,GB  & \multicolumn{2}{c}{out of budget} \\
\daWUK  & $(3,4$--$7)$ & 7.5\,s  & 0.4\,GB  & \multicolumn{2}{c}{out of budget} \\
\daPKa  & $(4,5$--$8)$ & 81\,s   & 29\,GB   & \multicolumn{2}{c}{out of budget} \\
\daPWT  & $(4,5$--$8)$ & 22\,s   & 7.6\,GB  & \multicolumn{2}{c}{out of budget} \\
\daLDR  & $(4,5$--$8)$ & 18\,s   & 7.7\,GB  & \multicolumn{2}{c}{out of budget} \\
\daRAE  & $(4,5$--$8)$ & 0.6\,s  & 0.2\,GB  & 307\,s   & 28\,GB \\
\daPKb  & $(4,5$--$8)$ & 3.7\,s  & 1.4\,GB  & 1{,}350\,s & 156\,GB \\
\daNAS  & $(4,5$--$8)$ & 25\,s   & 8.5\,GB  & 407\,s   & 61\,GB \\
\daDBL  & $(5,6$--$9)$ & 28\,s   & 2.7\,GB  & 4{,}261\,s & 257\,GB \\
\daDBL  & $(4,5$--$8)$ & 9.3\,s  & 1.1\,GB  & 212\,s   & 16\,GB \\
\daAST  & $(4,5$--$8)$ & 90\,s   & 10\,GB   & 148\,s   & 9\,GB \\
\daHEP  & $(3,4$--$7)$ & 167\,s  & 15\,GB   & 138\,s   & 4\,GB \\
\daEPN  & $(3,4$--$6)$ & 474\,s  & 26\,GB   & 2{,}807\,s & 2\,GB \\
\daYTB  & $(3,4$--$6)$ & 105\,s  & 7.8\,GB  & 605\,s   & 3.9\,GB \\
\bottomrule
\end{tabular}
\end{table}

\expsection{Spectrum Memory}
Table~\ref{tab:main} also reports peak memory.
%
The memory advantage is largest exactly where the per-cell state is largest: \cnd{} needs $156$GB on \daPKb{} and $257$GB on \daDBL{} at $r{=}5$, while \specnd{} needs $1.4$GB and $2.7$GB, reductions of $111\times$ and $95\times$; on the five out-of-budget graphs, \cnd{} exceeds $300$GB on every cell while \specnd{} stays under $29$GB.
%
On the social graphs the situation reverses: their $r$-clique counts are small, so \cnd{} is light ($2$--$4$GB), while \specnd{}'s residue machinery uses more ($8$--$26$GB).
%
The memory win therefore lives at the frontier where per-cell decomposition becomes infeasible, not on graphs that already fit.

\expsection{Effect of the Chain Certificate}
The certificate's power is seen most sharply by comparing a certified cell against recomputing the same cell with our own engine, \percell.
%
On \daDBL{} at $r{=}4$, the cell $s{=}8$ costs \percell{} $655$ seconds and \specnd{} $0.07$ seconds, a reduction of $9{,}364$ times; on \daDBL{} at $r{=}5$, the cell $s{=}9$ costs \percell{} over two hours and \specnd{} $0.15$ seconds.
%
On \daHEP{} and \daAST, the deepest cells ($(3,7)$ and $(4,8)$) exceed the two-hour per-cell budget for \percell, while \specnd{} produces them in $10$ seconds each.
%
This experiment demonstrates that the certificate, not merely the shared build, is what collapses the upper cells: once a pattern certifies, its remaining spectrum is arithmetic.

\expsection{Index Size and Query Time}
Table~\ref{tab:index} reports the size of \nsi{} and its query latency.
%
\nsi{} stores a whole spectrum at $0.045$ to $40$ bytes per $r$-clique: on \daWIT, the four-cell spectrum of $338$ million $r$-cliques fits in $14.5$MB, and on \daDBL{} at $r{=}5$ the spectrum of $263$ million $r$-cliques fits in $96$MB.
%
Point queries take $81$ to $412$ nanoseconds and full-spectrum queries $154$ to $678$ nanoseconds, using a randomly sampled workload of $200$ thousand $r$-cliques per graph.
%
The index is built as a byproduct of the sweep at no measurable extra cost.
%
On social graphs the residue dictionaries dominate, so the bytes-per-$r$-clique figure rises to $34$--$40$, but the whole spectrum still occupies tens of megabytes.

% -------- Index table --------
\begin{table}[t]
\centering
\caption{Index size and query latency.}
\label{tab:index}
\small
\begin{tabular}{lrrrr}
\toprule
Graph & size & B / $r$-clique & point & spectrum \\
\midrule
\daWIT  & 14.5\,MB  & 0.045 & 81\,ns  & 154\,ns \\
\daDBL{} $r{=}5$ & 95.7\,MB & 0.38 & 130\,ns & 253\,ns \\
\daDBL{} $r{=}4$ & 34.5\,MB & 2.2  & 150\,ns & 275\,ns \\
\daHEP  & 33.5\,MB  & 10.5  & 200\,ns & 354\,ns \\
\daAST  & 142\,MB   & 15.6  & 371\,ns & 574\,ns \\
\daYTB  & 83.7\,MB  & 34.2  & 215\,ns & 328\,ns \\
\daEPN  & 59.2\,MB  & 40.3  & 299\,ns & 421\,ns \\
\bottomrule
\end{tabular}
\end{table}

\expsection{Certification Coverage}
The wins of the previous experiments track one quantity: the fraction of $r$-cliques certified per cell.
%
Table~\ref{tab:cert} reports it.
%
On web and matrix graphs, every cell above the boundary certifies $100\%$ of $r$-cliques, and on \daWUK{} every region is mergeable, so the whole spectrum is closed form with no peeling at all.
%
On collaboration graphs certification reaches $99.97\%$, leaving a residue of a few thousand patterns out of millions.
%
On social graphs, where cliques overlap loosely, certification falls to $79\%$ on \daYTB{} and $47\%$ on \daEPN, and the uncertified residue must be peeled; this is exactly where the time and memory advantages narrow.
%
This experiment explains the earlier results: certification coverage is the single variable that separates the order-of-magnitude wins from the modest ones.

% -------- Certification table --------
\begin{table}[t]
\centering
\caption{Certified $r$-cliques per cell above the boundary.}
\label{tab:cert}
\small
\begin{tabular}{llr}
\toprule
Domain & Graph & certified \\
\midrule
Web / Matrix & \daWIT, \daPKa, \daRAE, \dots & $100\%$ \\
Collaboration & \daDBL & $100\%$ \\
Collaboration & \daAST, \daHEP & $99.97\%$ \\
Social & \daYTB & $79$--$81\%$ \\
Social & \daEPN & $47$--$51\%$ \\
\bottomrule
\end{tabular}
\end{table}

\expsection{Effect of the Shared Build}
\specndbase{} shares the build but peels every cell in full.
%
\todo{E7: fill SpecND vs SpecND-Base on hepph/astro/epin from tods2:/home/wenqianz/nsi\_e7/ when it lands; isolates certificate vs shared build}

\expsection{The Diagonal Transfer}
Theorem~\ref{thm:diagonal} predicts that the core value at cell $(r{+}1,r{+}2)$ is bounded by the minimum over $r$-subcliques at cell $(r,r{+}1)$ minus one.
%
We tested the bound on all $4$-cliques of three graphs by comparing the exact cells $(3,4)$ and $(4,5)$: across $156$ million $4$-cliques the bound holds with zero violations, and it is an equality for $100\%$ of the $4$-cliques on \daGRQ{} and \daHEP{} and $23.4\%$ on \daEPN.
%
The equality on collaboration graphs shows that the diagonal recurrence holds there for every single $4$-clique, so a one-pass diagonal is information-theoretically real on such graphs.
